% ============================================================
% Chapter 5: Development
%   Section: The Labelling Tool
% Included from chapters/05-development/development.tex
% Source: labeller/README.md; labeller/app.py.
% ============================================================

\section{The Labelling Tool}
\label{sec:dev-labeller}

The custom validation set (Section~\ref{sec:custom-dataset}) is labelled
with a tool built for it: one Python file and one HTML page, using only
the standard library, so that labelling needs nothing installed and
none of the system's services running. It listens on the local machine
only and accepts writes only as JSON, so another website open in the
same browser cannot post facts into the repository through it.

\paragraph{One file per fact.} Every saved fact is its own JSON file,
numbered in the order of saving, with X-Fact's keys in X-Fact's order
followed by the custom set's own. Every hand-verified fact therefore
appears in the repository, and in a diff, on its own; a test holds the
shared keys to the X-Fact files so that the harness scores both sets
with the same code. When the set is complete, a \texttt{join} command
re-validates every file and merges them into one JSON Lines file; it
refuses, and writes nothing, if any file breaks a rule.

\paragraph{The daily batch.} The \emph{Today} tab asks the backend for
about ten articles drawn at random, seeded by the date, from what the
sources published in the last 30 days: one per source, languages
alternated, never one already labelled or proposed, none that the
mission screen leaves out, preferring the topic groups the balance
table is short of. Each comes with the one to three claims the
pipeline's own claim selector would check. \emph{Label} fills the form
from the article and claim; \emph{Skip} records why the claim is not
worth labelling --- an opinion, a prediction, a trivial statement, a
fragment, a duplicate --- and the share labelled is the claim
selector's precision.

\paragraph{The rules on save.} The tie-break rules that can be checked
mechanically are refused on save: a verdict other than
\texttt{UNVERIFIED} needs an evidence link; a claim backed only by its
protagonist's own source must be \texttt{UNVERIFIED}; evidence dated
after the article is refused. A verdict-by-group table shows the 25
cells against their targets while labelling.

\paragraph{The blind review.} The \emph{Review} tab shows 20\% of the
facts, chosen by a hash of their file names, without their first label,
note or own-source flag, and reports observed agreement and Cohen's
$\kappa$ on the first label, so that correcting a fact afterwards
cannot inflate it.

\paragraph{Synthetic refuted claims.} A \texttt{mutate} command builds
a \texttt{FALSE} fact from a natural \texttt{TRUE} one by substituting a
number, a date, an entity or a relation, keeps the source's article,
date and evidence, and records the provenance; the tool lists any
mutation whose source has since been relabelled
(Subsection~\ref{subsec:custom-mutations}).

The tool has 32 tests, run by the project's check, which write to a
temporary folder and never to the real set.
